\chapter{Playing the Game}

\begin{multicols}{2}

Game is played with dice, using d10 and d6, or ten sided dice and six sided dice

Every time a math operation results in a fractional number, round the number down.

\section{How skills work}

Skills are things a character can do and improve upon. To measure how good someone is, a skill value is calculated.

\textbf{Skill test:} A skill test consists of a dice roll to determine the degree of success in performing an action. Tests are made by rolling a 10-sided die and adding the result to the skill value. This sum is compared to a difficulty level to determine the degree of success of the action.

\textbf{Difficulty level (DL):} A level that defines how hard it is to overcome something. It is the number against which skill tests are made to determine the degree of success.

\textbf{Degrees of success:} Critical, hit, graze, and miss are the four degrees of success. Critical is the best success, and miss is the worst failure. Regardless of the type of skill test, a score that exceeds a DL or an opponent defense (defense skill) by 10 or more is a critical success, 5 or more is a hit, 0 or more is a graze, and anything less than that is a miss.

\textbf{Exploding die:} In skill tests, if the result of the d10 is 10, roll 1d6 and add that to the score, continuing to add while the result is 6. If the result of the d10 is 1, do not add anything and then roll 1d6 and subtract the result instead of adding. Continue subtracting while rolling 6.

\textbf{Safe and risky tests:} It is possible to perform some tests safely or riskily. In the safe form, roll 2 d10 and take the result closest to 5. Similarly, it is possible to perform tests in a risky manner. In this case, roll 2 d10 and take the result farthest from 5. On a draw, use the average value or roll again. This mechanic can only be used when indicated in the skill or ability or when the GM allows it.

\textbf{On defining DLs:} the question to ask when estimating a value for DLs is "what skill level should have equal odds of success and failure?". Don't think about the number to be matched. Skill values exist usually between 0 and 10. Penalties or bonuses may push them to -5 or 15, but almost never beyond that. Consider that a DL of -5 is trivial, 0 is average, 5 is difficult, and 10 is very difficult to a human.

\textbf{Using tests:} A test happens whenever there is some degree of uncertainty over an outcome. If the scene makes a test trivial or impossible, do not call for a test and simply concede success or communicate failure. 

The more specific the player is about what their character is doing, the less a test becomes necessary. Simply concede that what they are trying to do works. Sometimes, a player may describe a correct course of action for their character because of their knowledge as a player. If that action is suspected to be beyond the character's abilities, the advice is to ask for a skill test to determine if the character would really do or know that. Conversely, often the character is smarter than the player, so the player can ask the character if they would know something through a test.

\textbf{Repeating tests:} tests are meant to change game state. That means that the character will have to do something else after the test, or some character value will be updated. If no character value was updated and nothing about the scene changed, do not call for a new test. Either use the same success from the previous test or use skill value against the new DL. Each specific section where that should happen has some advice on how to handle that.

\textbf{Group tests:} The skill value used for a group follows one of the following patterns: use the best value of the team (like in a detection test), use worst value of the team (like the slowest character is caught first) or use the best value but add bonus per helper (like in a grapple test). There will be instructions on which to use for each situation.

\section{Game loops}

The game runs on four loops: the big loop frames the campaign, scenes resolve action moment-to-moment, exploration moves the party through the world, and downtime grows characters between ventures.

\subsection{Big loop}

This loop depends on what kind of game is desired. Some games can be dungeon crawlers or heists, where life is cheap and characters are just a set of numbers used to beat a challenge. Others can be arena simulators, where characters are pitted against each other and fight like in a war game. Other games are stories, where characters have names, backstories and objectives. 

This system has tools to enable any of those kinds of games, or all at the same time. The exploration system for dungeons and survival, the combat system for battles, and the social system for storytelling and characterization.

\subsection{Scene:} This is where players describe their character's actions most precisely. It's where the action lives.

\textbf{Grid:} All game scenes where there is combat require a grid to be played. The other types can be estimated. Character positioning is very important and must be agreed upon by everyone at all times. There is a tight relationship between the action economy and movement, so approximating will, in most situations, devalue character build choices.

The grid can be either a hexgrid (recommended) or a grid of squares. Each space is, in general, equivalent to 1 action point in movement or 1m in diameter. Large creatures will take more spaces. 

\textbf{Modes:} There are four modes of scene: combat, chase, stealth and social. The types of scene can bleed into each other. One character may be trying to talk or sneak while others are trying to fight. Each kind of mode has its strategy to combine with the others. 

\subsubsection{Combat:} This is where most of the strategic element of the game lives. Players will use character skills, gear, abilities and the game's pre-defined actions to beat challenges. It is the part where personal opinions and GM adjudications influence outcomes the least. The main enjoyment of this should come from beating a challenge. For that to be true, the rules need to be well defined and everyone should agree to what they are, much like a puzzle or chess game. 

Combats are played on a grid in rounds that represent approximately 3 seconds of in game time, where each character acts in turns, which are decided on a "first to ask, first to play" basis that can be contested by other characters. Every action has a cost in time, represented by AP and often a cost in stamina. The combat loop revolves around draining opponent time and energy resources while preserving your own to have better chances of causing injuries while preventing damage to yourself through active defense.

There was special care put into making weapons and armor that would play like real medieval gear did. The same can be said for character's characteristics, like size, strength and stamina. Plate armor deflects almost any weapon attack. 

The game provides a list of disengagement options that prevent every combat from ending in death and diversify strategies and value different kinds of characters. They rely on the morale system with its limit actions and also on chase scenes. These options act like bridges between combat and the rest of the systems in the game.

\begin{proplist}
  \item \textbf{Negotiation:} Happens at the morale test. Can lead to capitulation or a truce.
  \item \textbf{Escape:} one side decides to run away and a chase scene is started.
  \item \textbf{Capitulation:} Each individual character stands down before dying. Combat ends when a round starts with all the team down. Bleeding punishes characters who continue fighting, so giving up early can be smart.
\end{proplist}

Not every combat must have defeating an enemy as the main objective. This game has mechanics to cover several different objectives for a combat scene. Here are a few:

\begin{proplist}
  \item \textbf{Defeat:} The point is to defeat all enemies and ends when one team dies, surrenders or flees.
  \item \textbf{Defend:} Ends when time runs out and the target is protected.
  \item \textbf{Defeat Target:} A specific target must be eliminated. One can flee after defeating the main target.
  \item \textbf{Escort/Kidnapping:} The objective is to defend one character and trigger a chase with it or eliminate all offenders.
\end{proplist}

\subsubsection{Chase}

Chases include fighting on the run and escaping. They are a major tool both as an alternative to combats to the death and as a repositioning strategy. 

A chase happens in distance tiers, where groups are within a given distance from each other. Different actions may become available depending on how far the groups are.

Characters have Stamina, gear burden and movement speeds, which come together into deciding who outruns who. Parties can decide to leave their slower members behind to be able to improve their speed, sacrificing someone to run away or just to make sure to catch enemies, if the advantage is overwhelming.

Chases can initiate combat again in a different location. The team who is escaping has the cards in their hands to decide where they are going. They can pick a location where they have an advantage. The system provides co-creation tools, where the player can try to pull the strings to find advantageous terrain.

Here are the predicted objectives for a chase scene:

\begin{proplist}
  \item \textbf{Reaching safety:} Finding a strong defensive point or allies is a form of escape.
  \item \textbf{Hiding:} The objective is to get enemies to lose track of the team.
  \item \textbf{Outrunning:} The objective is to convince enemies that they can't catch the team.
  \item \textbf{Separating:} The objective is to force opponents to split the party and risk being picked off.
  \item \textbf{Catching:} The objective is to pin down one or more characters of the opposite team and restart combat.
  \item \textbf{Fight on the run:} The objective is to stay in the chase mode and win the fight in it.
\end{proplist}

\subsubsection{Social} 

This is where players that enjoy acting will find themselves most comfortable. It involves players and the GM interpreting their characters to each other. The main role of this system is to structure the conversation, making characters commit to decisions and to their own characteristics.

The idea behind social interactions is to let the acting happen. The rules are intentionally light. Adding too much system on top of acting shuts it down instead of enhancing it. However, this is not meant to be free form conversation. Character skill must be taken into consideration for social skill to be a characteristic. The mechanics proposed in this book, such as the skill test with Interpretation bonus and the Contacts are designed to do just that.

Despite that, not everyone needs to be an actor. Saying something like: "I try to convince him to give me the bag" is valid, but doing so will forego the chance to convince the GM to give a skill bonus. This means that not acting will tend to result in a lower chance of success. 

Common kinds of social actions are:

\begin{proplist}
  \item \textbf{Performance:} Playing a role and pretending to be someone else.
  \item \textbf{Persuasion:} Trying to convince someone to believe you or doing something for you.
  \item \textbf{Negotiation:} Trying to get the best bargain out of a trade.
  \item \textbf{Seduction:} Trying to make someone like you personally.
  \item \textbf{Intimidation:} Coercing through threats.
\end{proplist}

For each of these actions, there should be irreversible consequences that drive the story forward. If there are no such consequences, the test can just be repeated indefinitely.

\subsubsection{Stealth}

Stealth scenes rely both on player descriptions and GM decisions. Each stealth skill test happens in the context of a stealth event, a description from the player of the character's action, which will help determine the difficulty of the skill test to be made. The player is encouraged to ask questions about what can be used, especially questions like "Do I see an object that I could hide behind?" or "Are there any bushes where I can camouflage myself in?". This kind of specific question is much easier to be accepted by the GM, since it unloads the cognitive burden from them. The players should take advantage of that fact and try to suggest things that are for their advantage. Nevertheless, the GM's word is final. If players try to force the issue and get easy and too convenient props all the time, the GM should put limits on them.

It is very unwieldy to attempt to model every aspect that will influence the difficulty of a stealth event. It would require tables upon tables of modifiers, which would slow the game down excessively. The real factors that make the test easier or harder are very situational and they can be expanded to excruciating detail. That is why this game offers a list of stealth modifiers as guidelines, not rules. That said, nothing stops someone from making stealth into a puzzle with pre-defined test values. That is a burden the GM will inflict on himself through game preparation and the game has the tables to help achieve that experience.

The intent is to achieve a middle ground between a simulation and a fluid game. The scenario co-creation tools are there to present hooks for the players to help themselves succeed while they help the GM.

A stealth scene can be done with or without the grid. It usually happens in the grid when there's risk of combat. A failure in stealth may trigger the beginning of any of the other scene modes.

Stealth scenes usually have three kinds of objectives:

\textbf{Stealth:} The objective is to go unnoticed.

\textbf{Theft:} The objective is to acquire an object. 

\textbf{Sabotage:} The objective is to interact with a specific object. 

\subsection{Exploration} 

A point crawl system, where players have a set of resources that they need to manage in order to navigate between points of interest. It is designed to be harsh, like real exploration is. Having survival gear, travel supplies, exploration skills and knowledge of the region are essential for success. Failing to do so will send players into real life threatening situations. 

A day is divided into 4 hour blocks, where each character can have independent action within the point crawl node. Each action costs energy in the form of hunger, thirst and exhaustion and each node can have multiple options of events to choose from. "Is it worth it for my character to make the effort when I am not skilled at something?", "I am very tired, but I am the only skilled character for this. Do I leave it to others or do I sacrifice my resources for an eventual fight?", "Is it worth spending time in this specific event or should I try to get something more useful?". These are the kinds of questions that players are meant to ask on a regular basis.

A common problem of this kind of system is that they tend to devolve into a slow and painful death march. This should be avoided, but without eliminating the threat of starvation or dehydration. The idea is to force players to make dramatic decisions and scenes to survive and reach their goal, and the events system is built to make those decisions happen. Scavengers, slavers, pilgrims, "friendly" monsters. What are characters willing to trade for their lives? 


\subsection{Downtime} 

This is where characters grow and prepare. It is where they improve their skills, gain new abilities, make new contacts, buy items, and do all the things that don't require a scene to be understood. Players will manage time, money and contacts to improve their characters through XP and karma.

\end{multicols}


